% 08_job_flow.tex -- one inference job from app to NPU and back: sequence, interrupt logic, DMA descriptors, coherence.
% Build: tectonic 08_job_flow.tex
\documentclass[tikz,border=10pt,12pt]{standalone}
% soc_style.tex -- shared look for the SoC figures (same palette as the slide decks and the
% core register placement figure). Every figure does % soc_style.tex -- shared look for the SoC figures (same palette as the slide decks and the
% core register placement figure). Every figure does % soc_style.tex -- shared look for the SoC figures (same palette as the slide decks and the
% core register placement figure). Every figure does \input{soc_style}.
\usepackage[T1]{fontenc}
\usepackage[scaled=0.95]{helvet}
\renewcommand{\familydefault}{\sfdefault}
\usepackage{booktabs}
\usepackage{amsmath}
\usepackage{amssymb}
\usetikzlibrary{arrows.meta,calc,positioning,fit,shapes.geometric,shapes.misc}

% ---------------- palette ----------------
\definecolor{pcb}{HTML}{0B2E33}
\definecolor{pcbmid}{HTML}{1F5C61}
\definecolor{tint}{HTML}{E6F0EF}
\definecolor{copper}{HTML}{B5651D}
\definecolor{coppert}{HTML}{F6E8DC}
\definecolor{ink}{HTML}{1F2A2E}
\definecolor{muted}{HTML}{5B6B70}
\definecolor{hardf}{HTML}{DCE6F2}
\definecolor{harde}{HTML}{2F4A7A}
\definecolor{good}{HTML}{2E7D32}
\definecolor{bad}{HTML}{B3261E}
% one colour per kind of block: border (d) / fill (f)
\definecolor{cpud}{HTML}{2B5C9E} \definecolor{cpuf}{HTML}{D6E4F5}
\definecolor{gpud}{HTML}{2E7D32} \definecolor{gpuf}{HTML}{D5EDD0}
\definecolor{npud}{HTML}{0B2E33} \definecolor{npuf}{HTML}{BFE0DB}
\definecolor{memd}{HTML}{5E3A8C} \definecolor{memf}{HTML}{E4DAF3}
\definecolor{busd}{HTML}{B5651D} \definecolor{busf}{HTML}{F6E8DC}
\definecolor{prd}{HTML}{546E7A}  \definecolor{prf}{HTML}{E3EAED}
\definecolor{secd}{HTML}{B3261E} \definecolor{secf}{HTML}{F7D9D6}
\definecolor{pwrd}{HTML}{8A6800} \definecolor{pwrf}{HTML}{FFF0B8}
\definecolor{medd}{HTML}{8E2463} \definecolor{medf}{HTML}{F5D6E6}
\definecolor{skd}{HTML}{2E7D32}  \definecolor{skf}{HTML}{BFE3B4}   % skew buffers, as in the register figure

\tikzset{
  blk/.style 2 args={draw=#1, fill=#2, rounded corners=4pt, thick, align=center,
                     font=\scriptsize, inner sep=3pt},
  cpu/.style={blk={cpud}{cpuf}},  gpu/.style={blk={gpud}{gpuf}},
  npu/.style={blk={npud}{npuf}},  mem/.style={blk={memd}{memf}},
  bus/.style={blk={busd}{busf}},  per/.style={blk={prd}{prf}},
  sec/.style={blk={secd}{secf}},  pwr/.style={blk={pwrd}{pwrf}},
  med/.style={blk={medd}{medf}},
  comb/.style ={draw=copper, fill=coppert, rounded corners=2pt, thick, align=center,
                font=\scriptsize, inner sep=2pt},                       % combinational logic
  hard/.style ={draw=harde, fill=hardf, thick, dashed, align=center,
                font=\scriptsize, inner sep=2pt},                       % hard macro / off-chip
  sram/.style ={draw=memd, fill=memf, thick, align=center, font=\scriptsize,
                inner sep=2pt},                                         % SRAM array, not flip-flops
  ff/.style   ={draw=#1, fill=#1!30, line width=0.35pt},                % one flip-flop square
  ffgone/.style={draw=muted, fill=white, densely dashed, line width=0.35pt},
  panel/.style={draw=muted, rounded corners=6pt, line width=0.9pt},
  ptitle/.style={anchor=north west, font=\normalsize\bfseries, text=ink, inner sep=0pt},
  note/.style ={font=\scriptsize, text=ink, align=left, inner sep=0pt},
  lab/.style  ={font=\tiny, text=muted, inner sep=1pt},
  axi/.style  ={{Stealth[length=2.6mm]}-{Stealth[length=2.6mm]}, line width=2.2pt, copper},
  axi1/.style ={-{Stealth[length=2.6mm]}, line width=2.2pt, copper},
  apb/.style  ={{Stealth[length=1.8mm]}-{Stealth[length=1.8mm]}, line width=1pt, copper},
  apb1/.style ={-{Stealth[length=1.8mm]}, line width=1pt, copper},
  dat/.style  ={-{Stealth[length=1.8mm]}, thick, pcbmid},
  ctl/.style  ={-{Stealth[length=1.8mm]}, densely dashed, thick, muted},
  irq/.style  ={-{Stealth[length=2mm]}, densely dotted, line width=1.3pt, bad},
  num/.style  ={circle, fill=copper, text=white, font=\scriptsize\bfseries, inner sep=1pt,
                minimum size=0.46cm},
  ours/.style ={draw=bad, line width=1.2pt, densely dashed, rounded corners=4pt},
}

% ---------------- helpers ----------------
\newcommand{\rn}[1]{\texttt{#1}}
% \ffrow{x}{y}{n}{cell}{style}: n flip-flops in a row, bottom-left corner at (x,y)
\newcommand{\ffrow}[5]{\foreach \q in {1,...,#3}{\draw[#5] ({#1+(\q-1)*#4},{#2}) rectangle ++(#4,#4);}}
% \ffcol{x}{y}{n}{cell}{style}: n flip-flops in a column, top-left corner at (x,y)
\newcommand{\ffcol}[5]{\foreach \q in {1,...,#3}{\draw[#5] ({#1},{#2-\q*#4}) rectangle ++(#4,#4);}}
% \ffgrid{x}{ytop}{n}{per row}{cell}{style}: n flip-flops, filled row by row from the top-left
\newcommand{\ffgrid}[6]{\pgfmathtruncatemacro{\gl}{#3-1}%
  \foreach \b in {0,...,\gl}{\pgfmathtruncatemacro{\gc}{mod(\b,#4)}\pgfmathtruncatemacro{\gr}{div(\b,#4)}%
    \draw[#6] ({#1+\gc*#5},{#2-(\gr+1)*#5}) rectangle ++(#5,#5);}}

% every figure: same page size, title, footer
\newcommand{\canvas}{\useasboundingbox (-0.3,-0.2) rectangle (38.6,30.4);}
\newcommand{\figtitle}[3]{%
  \node[num, minimum size=0.9cm, font=\Large\bfseries] at (0.45,29.55) {#1};
  \node[anchor=west, font=\LARGE\bfseries, text=pcb] at (1.15,29.55) {#2};
  \node[anchor=west, font=\small, text=ink] at (0,28.7) {#3};}
\newcommand{\figfoot}[1]{\node[anchor=west, font=\scriptsize, text=muted] at (0,0.25) {#1};}
% \panel{x0}{y0}{x1}{y1}{title}
\newcommand{\panel}[5]{\draw[panel] (#1,#2) rectangle (#3,#4);
  \node[ptitle] at ({#1+0.3},{#4-0.25}) {#5};}

% ---------------- timing diagrams: 1 cycle = \cw cm, low at y, high at y+0.45 ----------------
\newcommand{\cw}{1.0}
\newcommand{\wclk}[3]{% x0 y ncycles
  \foreach \k in {1,...,#3}{\draw[pcb, thick] ({#1+(\k-1)*\cw},#2) -- ++(0,0.45) -- ++(0.5*\cw,0)
     -- ++(0,-0.45) -- ++(0.5*\cw,0);}}
\newcommand{\wsig}[3]{% x0 y {level per cycle, e.g. 0,1,1,0}
  \gdef\wprev{x}%
  \foreach \v [count=\k from 0] in {#3}{%
    \draw[pcb, thick] ({#1+\k*\cw},{#2+0.45*\v}) -- ++(\cw,0);
    \if x\wprev\else\ifnum\wprev=\v\else\draw[pcb, thick] ({#1+\k*\cw},#2) -- ++(0,0.45);\fi\fi
    \global\let\wprev\v}}
\newcommand{\wbus}[5]{% x0 y c0 c1 text: one value held from cycle c0 to cycle c1
  \draw[pcbmid, thick, fill=white] ({#1+(#3)*\cw},{#2+0.225}) -- ({#1+(#3)*\cw+0.08},{#2+0.45})
    -- ({#1+(#4)*\cw-0.08},{#2+0.45}) -- ({#1+(#4)*\cw},{#2+0.225}) -- ({#1+(#4)*\cw-0.08},#2)
    -- ({#1+(#3)*\cw+0.08},#2) -- cycle;
  \node[font=\tiny\ttfamily] at ({#1+((#3)+(#4))*\cw/2},{#2+0.225}) {#5};}
\newcommand{\widle}[4]{% x0 y c0 c1: bus idle (no value)
  \draw[muted, thick] ({#1+(#3)*\cw},{#2+0.225}) -- ({#1+(#4)*\cw},{#2+0.225});}
\newcommand{\wname}[3]{\node[anchor=east, font=\scriptsize\ttfamily] at ({#1-0.15},{#2+0.225}) {#3};}
.
\usepackage[T1]{fontenc}
\usepackage[scaled=0.95]{helvet}
\renewcommand{\familydefault}{\sfdefault}
\usepackage{booktabs}
\usepackage{amsmath}
\usepackage{amssymb}
\usetikzlibrary{arrows.meta,calc,positioning,fit,shapes.geometric,shapes.misc}

% ---------------- palette ----------------
\definecolor{pcb}{HTML}{0B2E33}
\definecolor{pcbmid}{HTML}{1F5C61}
\definecolor{tint}{HTML}{E6F0EF}
\definecolor{copper}{HTML}{B5651D}
\definecolor{coppert}{HTML}{F6E8DC}
\definecolor{ink}{HTML}{1F2A2E}
\definecolor{muted}{HTML}{5B6B70}
\definecolor{hardf}{HTML}{DCE6F2}
\definecolor{harde}{HTML}{2F4A7A}
\definecolor{good}{HTML}{2E7D32}
\definecolor{bad}{HTML}{B3261E}
% one colour per kind of block: border (d) / fill (f)
\definecolor{cpud}{HTML}{2B5C9E} \definecolor{cpuf}{HTML}{D6E4F5}
\definecolor{gpud}{HTML}{2E7D32} \definecolor{gpuf}{HTML}{D5EDD0}
\definecolor{npud}{HTML}{0B2E33} \definecolor{npuf}{HTML}{BFE0DB}
\definecolor{memd}{HTML}{5E3A8C} \definecolor{memf}{HTML}{E4DAF3}
\definecolor{busd}{HTML}{B5651D} \definecolor{busf}{HTML}{F6E8DC}
\definecolor{prd}{HTML}{546E7A}  \definecolor{prf}{HTML}{E3EAED}
\definecolor{secd}{HTML}{B3261E} \definecolor{secf}{HTML}{F7D9D6}
\definecolor{pwrd}{HTML}{8A6800} \definecolor{pwrf}{HTML}{FFF0B8}
\definecolor{medd}{HTML}{8E2463} \definecolor{medf}{HTML}{F5D6E6}
\definecolor{skd}{HTML}{2E7D32}  \definecolor{skf}{HTML}{BFE3B4}   % skew buffers, as in the register figure

\tikzset{
  blk/.style 2 args={draw=#1, fill=#2, rounded corners=4pt, thick, align=center,
                     font=\scriptsize, inner sep=3pt},
  cpu/.style={blk={cpud}{cpuf}},  gpu/.style={blk={gpud}{gpuf}},
  npu/.style={blk={npud}{npuf}},  mem/.style={blk={memd}{memf}},
  bus/.style={blk={busd}{busf}},  per/.style={blk={prd}{prf}},
  sec/.style={blk={secd}{secf}},  pwr/.style={blk={pwrd}{pwrf}},
  med/.style={blk={medd}{medf}},
  comb/.style ={draw=copper, fill=coppert, rounded corners=2pt, thick, align=center,
                font=\scriptsize, inner sep=2pt},                       % combinational logic
  hard/.style ={draw=harde, fill=hardf, thick, dashed, align=center,
                font=\scriptsize, inner sep=2pt},                       % hard macro / off-chip
  sram/.style ={draw=memd, fill=memf, thick, align=center, font=\scriptsize,
                inner sep=2pt},                                         % SRAM array, not flip-flops
  ff/.style   ={draw=#1, fill=#1!30, line width=0.35pt},                % one flip-flop square
  ffgone/.style={draw=muted, fill=white, densely dashed, line width=0.35pt},
  panel/.style={draw=muted, rounded corners=6pt, line width=0.9pt},
  ptitle/.style={anchor=north west, font=\normalsize\bfseries, text=ink, inner sep=0pt},
  note/.style ={font=\scriptsize, text=ink, align=left, inner sep=0pt},
  lab/.style  ={font=\tiny, text=muted, inner sep=1pt},
  axi/.style  ={{Stealth[length=2.6mm]}-{Stealth[length=2.6mm]}, line width=2.2pt, copper},
  axi1/.style ={-{Stealth[length=2.6mm]}, line width=2.2pt, copper},
  apb/.style  ={{Stealth[length=1.8mm]}-{Stealth[length=1.8mm]}, line width=1pt, copper},
  apb1/.style ={-{Stealth[length=1.8mm]}, line width=1pt, copper},
  dat/.style  ={-{Stealth[length=1.8mm]}, thick, pcbmid},
  ctl/.style  ={-{Stealth[length=1.8mm]}, densely dashed, thick, muted},
  irq/.style  ={-{Stealth[length=2mm]}, densely dotted, line width=1.3pt, bad},
  num/.style  ={circle, fill=copper, text=white, font=\scriptsize\bfseries, inner sep=1pt,
                minimum size=0.46cm},
  ours/.style ={draw=bad, line width=1.2pt, densely dashed, rounded corners=4pt},
}

% ---------------- helpers ----------------
\newcommand{\rn}[1]{\texttt{#1}}
% \ffrow{x}{y}{n}{cell}{style}: n flip-flops in a row, bottom-left corner at (x,y)
\newcommand{\ffrow}[5]{\foreach \q in {1,...,#3}{\draw[#5] ({#1+(\q-1)*#4},{#2}) rectangle ++(#4,#4);}}
% \ffcol{x}{y}{n}{cell}{style}: n flip-flops in a column, top-left corner at (x,y)
\newcommand{\ffcol}[5]{\foreach \q in {1,...,#3}{\draw[#5] ({#1},{#2-\q*#4}) rectangle ++(#4,#4);}}
% \ffgrid{x}{ytop}{n}{per row}{cell}{style}: n flip-flops, filled row by row from the top-left
\newcommand{\ffgrid}[6]{\pgfmathtruncatemacro{\gl}{#3-1}%
  \foreach \b in {0,...,\gl}{\pgfmathtruncatemacro{\gc}{mod(\b,#4)}\pgfmathtruncatemacro{\gr}{div(\b,#4)}%
    \draw[#6] ({#1+\gc*#5},{#2-(\gr+1)*#5}) rectangle ++(#5,#5);}}

% every figure: same page size, title, footer
\newcommand{\canvas}{\useasboundingbox (-0.3,-0.2) rectangle (38.6,30.4);}
\newcommand{\figtitle}[3]{%
  \node[num, minimum size=0.9cm, font=\Large\bfseries] at (0.45,29.55) {#1};
  \node[anchor=west, font=\LARGE\bfseries, text=pcb] at (1.15,29.55) {#2};
  \node[anchor=west, font=\small, text=ink] at (0,28.7) {#3};}
\newcommand{\figfoot}[1]{\node[anchor=west, font=\scriptsize, text=muted] at (0,0.25) {#1};}
% \panel{x0}{y0}{x1}{y1}{title}
\newcommand{\panel}[5]{\draw[panel] (#1,#2) rectangle (#3,#4);
  \node[ptitle] at ({#1+0.3},{#4-0.25}) {#5};}

% ---------------- timing diagrams: 1 cycle = \cw cm, low at y, high at y+0.45 ----------------
\newcommand{\cw}{1.0}
\newcommand{\wclk}[3]{% x0 y ncycles
  \foreach \k in {1,...,#3}{\draw[pcb, thick] ({#1+(\k-1)*\cw},#2) -- ++(0,0.45) -- ++(0.5*\cw,0)
     -- ++(0,-0.45) -- ++(0.5*\cw,0);}}
\newcommand{\wsig}[3]{% x0 y {level per cycle, e.g. 0,1,1,0}
  \gdef\wprev{x}%
  \foreach \v [count=\k from 0] in {#3}{%
    \draw[pcb, thick] ({#1+\k*\cw},{#2+0.45*\v}) -- ++(\cw,0);
    \if x\wprev\else\ifnum\wprev=\v\else\draw[pcb, thick] ({#1+\k*\cw},#2) -- ++(0,0.45);\fi\fi
    \global\let\wprev\v}}
\newcommand{\wbus}[5]{% x0 y c0 c1 text: one value held from cycle c0 to cycle c1
  \draw[pcbmid, thick, fill=white] ({#1+(#3)*\cw},{#2+0.225}) -- ({#1+(#3)*\cw+0.08},{#2+0.45})
    -- ({#1+(#4)*\cw-0.08},{#2+0.45}) -- ({#1+(#4)*\cw},{#2+0.225}) -- ({#1+(#4)*\cw-0.08},#2)
    -- ({#1+(#3)*\cw+0.08},#2) -- cycle;
  \node[font=\tiny\ttfamily] at ({#1+((#3)+(#4))*\cw/2},{#2+0.225}) {#5};}
\newcommand{\widle}[4]{% x0 y c0 c1: bus idle (no value)
  \draw[muted, thick] ({#1+(#3)*\cw},{#2+0.225}) -- ({#1+(#4)*\cw},{#2+0.225});}
\newcommand{\wname}[3]{\node[anchor=east, font=\scriptsize\ttfamily] at ({#1-0.15},{#2+0.225}) {#3};}
.
\usepackage[T1]{fontenc}
\usepackage[scaled=0.95]{helvet}
\renewcommand{\familydefault}{\sfdefault}
\usepackage{booktabs}
\usepackage{amsmath}
\usepackage{amssymb}
\usetikzlibrary{arrows.meta,calc,positioning,fit,shapes.geometric,shapes.misc}

% ---------------- palette ----------------
\definecolor{pcb}{HTML}{0B2E33}
\definecolor{pcbmid}{HTML}{1F5C61}
\definecolor{tint}{HTML}{E6F0EF}
\definecolor{copper}{HTML}{B5651D}
\definecolor{coppert}{HTML}{F6E8DC}
\definecolor{ink}{HTML}{1F2A2E}
\definecolor{muted}{HTML}{5B6B70}
\definecolor{hardf}{HTML}{DCE6F2}
\definecolor{harde}{HTML}{2F4A7A}
\definecolor{good}{HTML}{2E7D32}
\definecolor{bad}{HTML}{B3261E}
% one colour per kind of block: border (d) / fill (f)
\definecolor{cpud}{HTML}{2B5C9E} \definecolor{cpuf}{HTML}{D6E4F5}
\definecolor{gpud}{HTML}{2E7D32} \definecolor{gpuf}{HTML}{D5EDD0}
\definecolor{npud}{HTML}{0B2E33} \definecolor{npuf}{HTML}{BFE0DB}
\definecolor{memd}{HTML}{5E3A8C} \definecolor{memf}{HTML}{E4DAF3}
\definecolor{busd}{HTML}{B5651D} \definecolor{busf}{HTML}{F6E8DC}
\definecolor{prd}{HTML}{546E7A}  \definecolor{prf}{HTML}{E3EAED}
\definecolor{secd}{HTML}{B3261E} \definecolor{secf}{HTML}{F7D9D6}
\definecolor{pwrd}{HTML}{8A6800} \definecolor{pwrf}{HTML}{FFF0B8}
\definecolor{medd}{HTML}{8E2463} \definecolor{medf}{HTML}{F5D6E6}
\definecolor{skd}{HTML}{2E7D32}  \definecolor{skf}{HTML}{BFE3B4}   % skew buffers, as in the register figure

\tikzset{
  blk/.style 2 args={draw=#1, fill=#2, rounded corners=4pt, thick, align=center,
                     font=\scriptsize, inner sep=3pt},
  cpu/.style={blk={cpud}{cpuf}},  gpu/.style={blk={gpud}{gpuf}},
  npu/.style={blk={npud}{npuf}},  mem/.style={blk={memd}{memf}},
  bus/.style={blk={busd}{busf}},  per/.style={blk={prd}{prf}},
  sec/.style={blk={secd}{secf}},  pwr/.style={blk={pwrd}{pwrf}},
  med/.style={blk={medd}{medf}},
  comb/.style ={draw=copper, fill=coppert, rounded corners=2pt, thick, align=center,
                font=\scriptsize, inner sep=2pt},                       % combinational logic
  hard/.style ={draw=harde, fill=hardf, thick, dashed, align=center,
                font=\scriptsize, inner sep=2pt},                       % hard macro / off-chip
  sram/.style ={draw=memd, fill=memf, thick, align=center, font=\scriptsize,
                inner sep=2pt},                                         % SRAM array, not flip-flops
  ff/.style   ={draw=#1, fill=#1!30, line width=0.35pt},                % one flip-flop square
  ffgone/.style={draw=muted, fill=white, densely dashed, line width=0.35pt},
  panel/.style={draw=muted, rounded corners=6pt, line width=0.9pt},
  ptitle/.style={anchor=north west, font=\normalsize\bfseries, text=ink, inner sep=0pt},
  note/.style ={font=\scriptsize, text=ink, align=left, inner sep=0pt},
  lab/.style  ={font=\tiny, text=muted, inner sep=1pt},
  axi/.style  ={{Stealth[length=2.6mm]}-{Stealth[length=2.6mm]}, line width=2.2pt, copper},
  axi1/.style ={-{Stealth[length=2.6mm]}, line width=2.2pt, copper},
  apb/.style  ={{Stealth[length=1.8mm]}-{Stealth[length=1.8mm]}, line width=1pt, copper},
  apb1/.style ={-{Stealth[length=1.8mm]}, line width=1pt, copper},
  dat/.style  ={-{Stealth[length=1.8mm]}, thick, pcbmid},
  ctl/.style  ={-{Stealth[length=1.8mm]}, densely dashed, thick, muted},
  irq/.style  ={-{Stealth[length=2mm]}, densely dotted, line width=1.3pt, bad},
  num/.style  ={circle, fill=copper, text=white, font=\scriptsize\bfseries, inner sep=1pt,
                minimum size=0.46cm},
  ours/.style ={draw=bad, line width=1.2pt, densely dashed, rounded corners=4pt},
}

% ---------------- helpers ----------------
\newcommand{\rn}[1]{\texttt{#1}}
% \ffrow{x}{y}{n}{cell}{style}: n flip-flops in a row, bottom-left corner at (x,y)
\newcommand{\ffrow}[5]{\foreach \q in {1,...,#3}{\draw[#5] ({#1+(\q-1)*#4},{#2}) rectangle ++(#4,#4);}}
% \ffcol{x}{y}{n}{cell}{style}: n flip-flops in a column, top-left corner at (x,y)
\newcommand{\ffcol}[5]{\foreach \q in {1,...,#3}{\draw[#5] ({#1},{#2-\q*#4}) rectangle ++(#4,#4);}}
% \ffgrid{x}{ytop}{n}{per row}{cell}{style}: n flip-flops, filled row by row from the top-left
\newcommand{\ffgrid}[6]{\pgfmathtruncatemacro{\gl}{#3-1}%
  \foreach \b in {0,...,\gl}{\pgfmathtruncatemacro{\gc}{mod(\b,#4)}\pgfmathtruncatemacro{\gr}{div(\b,#4)}%
    \draw[#6] ({#1+\gc*#5},{#2-(\gr+1)*#5}) rectangle ++(#5,#5);}}

% every figure: same page size, title, footer
\newcommand{\canvas}{\useasboundingbox (-0.3,-0.2) rectangle (38.6,30.4);}
\newcommand{\figtitle}[3]{%
  \node[num, minimum size=0.9cm, font=\Large\bfseries] at (0.45,29.55) {#1};
  \node[anchor=west, font=\LARGE\bfseries, text=pcb] at (1.15,29.55) {#2};
  \node[anchor=west, font=\small, text=ink] at (0,28.7) {#3};}
\newcommand{\figfoot}[1]{\node[anchor=west, font=\scriptsize, text=muted] at (0,0.25) {#1};}
% \panel{x0}{y0}{x1}{y1}{title}
\newcommand{\panel}[5]{\draw[panel] (#1,#2) rectangle (#3,#4);
  \node[ptitle] at ({#1+0.3},{#4-0.25}) {#5};}

% ---------------- timing diagrams: 1 cycle = \cw cm, low at y, high at y+0.45 ----------------
\newcommand{\cw}{1.0}
\newcommand{\wclk}[3]{% x0 y ncycles
  \foreach \k in {1,...,#3}{\draw[pcb, thick] ({#1+(\k-1)*\cw},#2) -- ++(0,0.45) -- ++(0.5*\cw,0)
     -- ++(0,-0.45) -- ++(0.5*\cw,0);}}
\newcommand{\wsig}[3]{% x0 y {level per cycle, e.g. 0,1,1,0}
  \gdef\wprev{x}%
  \foreach \v [count=\k from 0] in {#3}{%
    \draw[pcb, thick] ({#1+\k*\cw},{#2+0.45*\v}) -- ++(\cw,0);
    \if x\wprev\else\ifnum\wprev=\v\else\draw[pcb, thick] ({#1+\k*\cw},#2) -- ++(0,0.45);\fi\fi
    \global\let\wprev\v}}
\newcommand{\wbus}[5]{% x0 y c0 c1 text: one value held from cycle c0 to cycle c1
  \draw[pcbmid, thick, fill=white] ({#1+(#3)*\cw},{#2+0.225}) -- ({#1+(#3)*\cw+0.08},{#2+0.45})
    -- ({#1+(#4)*\cw-0.08},{#2+0.45}) -- ({#1+(#4)*\cw},{#2+0.225}) -- ({#1+(#4)*\cw-0.08},#2)
    -- ({#1+(#3)*\cw+0.08},#2) -- cycle;
  \node[font=\tiny\ttfamily] at ({#1+((#3)+(#4))*\cw/2},{#2+0.225}) {#5};}
\newcommand{\widle}[4]{% x0 y c0 c1: bus idle (no value)
  \draw[muted, thick] ({#1+(#3)*\cw},{#2+0.225}) -- ({#1+(#4)*\cw},{#2+0.225});}
\newcommand{\wname}[3]{\node[anchor=east, font=\scriptsize\ttfamily] at ({#1-0.15},{#2+0.225}) {#3};}

% \step{n}{y}{from x}{to x}{style}{label}
\newcommand{\step}[6]{%
  \draw[#5] (#3,#2) -- (#4,#2);
  \pgfmathsetmacro{\cx}{#3<#4 ? #3-0.35 : #3+0.35}%
  \node[num, minimum size=0.42cm, font=\tiny\bfseries] at (\cx,#2) {#1};
  \node[font=\tiny, anchor=south, fill=white, inner sep=0.5pt] at ({(#3+#4)/2},{#2+0.04}) {#6};}
\tikzset{lite/.style={-{Stealth[length=1.8mm]}, line width=1pt, copper},
         sw/.style={-{Stealth[length=1.8mm]}, thick, muted}}
\begin{document}
\begin{tikzpicture}[line join=round]
\canvas
\figtitle{8}{One AI job end to end: CPU, DRAM, NPU and the interrupt}
  {What happens between ``the app calls the model'' and ``the app has the answer''. Each arrow is one bus transfer,
   one wire or one software call, in time order from top to bottom.}

% =====================================================================
% (A) sequence diagram
% =====================================================================
\panel{0}{12.2}{38.3}{28.1}{Sequence of one job}
\foreach \x/\nm/\sty in {2.6/{app\\(CPU, user)}/cpu, 7.6/{driver\\(CPU, kernel)}/cpu, 12.6/{CPU caches}/cpu,
                         17.6/{DRAM}/mem, 22.6/{NPU registers}/npu, 27.6/{NPU DMA +\\command proc.}/npu,
                         32.3/{MAC array +\\post-processing}/npu, 36.6/{GIC}/per}{
  \node[\sty, minimum width=3.2cm, minimum height=0.9cm, font=\scriptsize\bfseries] at (\x,26.95) {\nm};
  \draw[muted, densely dashed] (\x,26.5) -- (\x,13.2);}
\draw[draw=npud, densely dashed, rounded corners=4pt, fill=npuf, fill opacity=0.35] (16.8,20.65) rectangle (33.2,18.35);
\node[font=\tiny\bfseries, text=npud, anchor=north east] at (33.15,20.62) {repeat 8--10 for every tile of every layer};
\step{1}{25.9}{2.6}{7.6}{sw}{run(model, input): system call}
\step{2}{25.1}{7.6}{12.6}{dat}{write input + command list (lines now dirty in cache)}
\step{3}{24.3}{12.6}{17.6}{dat}{cache clean: dirty lines $\rightarrow$ DRAM}
\step{4}{23.5}{7.6}{22.6}{lite}{AXI-Lite writes: CMD\_BASE = 0x8010\_0000, CMD\_COUNT = 12}
\step{5}{22.7}{7.6}{22.6}{lite}{AXI-Lite write: CTRL.START = 1 (``doorbell''); the driver now sleeps, CPU is free}
\step{6}{21.9}{22.6}{27.6}{dat}{start pulse}
\step{7}{21.1}{27.6}{17.6}{axi1}{AXI read burst: command list}
\step{8}{20.2}{27.6}{17.6}{axi1}{AXI read bursts: weights + input tile}
\step{9}{19.45}{27.6}{32.3}{dat}{tile $\rightarrow$ SRAM $\rightarrow$ feeders $\rightarrow$ array}
\step{10}{18.7}{32.3}{17.6}{axi1}{int8 results $\rightarrow$ DMA write bursts $\rightarrow$ DRAM}
\step{11}{17.8}{27.6}{22.6}{dat}{last tile: DONE = 1}
\step{12}{17.0}{22.6}{36.6}{irq}{IRQ wire goes high (IRQ\_STATUS AND IRQ\_MASK)}
\step{13}{16.2}{36.6}{7.6}{irq}{GIC interrupts a CPU core: it saves its PC and jumps to the driver's handler}
\step{14}{15.4}{7.6}{22.6}{lite}{read IRQ\_STATUS, write 1 to clear it: IRQ wire goes low}
\step{15}{14.6}{7.6}{12.6}{dat}{invalidate output lines}
\step{16}{13.8}{7.6}{2.6}{sw}{wake app}
\node[note, font=\tiny, anchor=north west, text width=4.4cm, fill=white, inner sep=2pt] at (33.6,25.6) {%
  \textbf{Legend}\\
  \textcolor{muted}{\textbf{grey}}: software call\\
  \textcolor{pcbmid}{\textbf{teal}}: inside one block\\
  \textcolor{copper}{\textbf{thin copper}}: AXI-Lite (registers)\\
  \textcolor{copper}{\textbf{thick copper}}: AXI bursts (data)\\
  \textcolor{bad}{\textbf{red dotted}}: interrupt wire};
\node[note, font=\tiny, anchor=north west, text width=4.4cm, fill=white, inner sep=2pt] at (33.6,22.6) {%
  \textbf{Your project today}: \rn{host.py} = steps 1--5 over UART (writes A, B, START), polls STATUS instead of 12--14,
  reads C over UART. No cache, no DMA, no IRQ.};

% =====================================================================
% (B) interrupt hardware
% =====================================================================
\panel{0}{0.7}{12.6}{11.9}{Interrupt hardware, gate by gate}
\node[font=\tiny, anchor=east] at (1.3,10.4) {done pulse};
\node[comb, minimum width=0.8cm, minimum height=0.8cm, font=\tiny\bfseries] (or1) at (2.0,10.2) {OR};
\node[comb, minimum width=0.8cm, minimum height=0.8cm, font=\tiny\bfseries] (and1) at (3.6,10.2) {AND};
\node[draw=npud, fill=npuf, thick, minimum width=1.0cm, minimum height=1.0cm, font=\tiny, align=center] (ff1) at (5.4,10.2) {D\quad Q\\[2pt]sticky};
\draw[npud, thick] (5.05,9.7) -- (5.2,9.83) -- (5.35,9.7);
\draw[dat] (1.3,10.4) -- (1.6,10.4);
\draw[dat] (or1) -- (and1); \draw[dat] (and1) -- (ff1);
\draw[dat] (ff1.east) -- (6.6,10.2) -- (6.6,9.3) -- (1.45,9.3) -- (1.45,10.0) -- (1.6,10.0);
\node[font=\tiny, anchor=north] at (3.6,9.65) {NOT(write 1 to clear)};
\draw[dat] (3.6,9.75) -- (3.6,9.8);
\node[lab, text=npud, font=\tiny\bfseries, anchor=south] at (5.4,10.75) {IRQ\_STATUS bit};
\node[comb, minimum width=0.8cm, minimum height=0.8cm, font=\tiny\bfseries] (and2) at (7.6,10.2) {AND};
\draw[dat] (6.6,10.2) -- (and2);
\node[draw=npud, fill=npuf, thick, minimum width=1.0cm, minimum height=0.6cm, font=\tiny] (msk) at (7.6,8.9) {MASK FF};
\draw[dat] (msk) -- (and2);
\node[comb, minimum width=0.8cm, minimum height=0.8cm, font=\tiny\bfseries] (or2) at (9.3,10.2) {OR};
\draw[dat] (and2) -- (or2);
\node[font=\tiny, anchor=north] at (9.3,9.7) {other events};
\draw[irq] (or2) -- (11.9,10.2);
\node[font=\tiny\bfseries, text=bad, anchor=south] at (10.9,10.25) {irq};
\node[note, font=\scriptsize, anchor=north west, text width=12.0cm] at (0.3,8.35) {%
  next Q = (Q OR done) AND NOT clear: it stays 1 until software clears it. Your \rn{mm\_regs} \rn{done\_flag} is the same circuit
  (cleared by START instead of a write-1-to-clear).};
\node[per, fill=white, minimum width=5.8cm, minimum height=2.5cm, text width=5.6cm, anchor=north west, font=\tiny] at (0.3,7.1)
  {\textbf{GIC distributor}\\for each of $\approx$300 IRQ inputs:\\enable bit $\cdot$ 8-bit priority $\cdot$ target core\\[2pt]
   picks the highest-priority pending IRQ};
\node[per, fill=white, minimum width=5.8cm, minimum height=2.5cm, text width=5.6cm, anchor=north west, font=\tiny] at (6.45,7.1)
  {\textbf{GIC CPU interface}\\raises the IRQ line of one core\\[2pt]core: finish the current instruction, save PC, jump to the vector, run the handler,
   then tell the GIC ``end of interrupt''};
\draw[irq] (6.1,5.85) -- (6.45,5.85);
\node[note, font=\scriptsize, anchor=north west, text width=12.0cm] at (0.3,4.4) {%
  \textbf{Polling vs interrupt}: polling reads STATUS again and again (your \rn{host.py}): simple, but it keeps the CPU busy.
  An interrupt lets the CPU sleep or do other work until the NPU calls. Interrupt cost: a few hundred cycles to save state and run the handler.\\[3pt]
  \textbf{Level vs edge}: this IRQ is level-sensitive: the wire stays high until software clears the bit, so it cannot be missed.};

% =====================================================================
% (C) DMA descriptors
% =====================================================================
\panel{12.9}{0.7}{25.4}{11.9}{DMA descriptors: a to-do list in DRAM}
\node[anchor=north west, font=\scriptsize\bfseries] at (13.2,11.15) {one descriptor = 40 bytes:};
\foreach \nm/\w [count=\k from 0] in {source addr/64, dest addr/64, length/32, src stride/32, dst stride/32, rows/16, ctrl/16, next/64}{
  \ifnum\k=0 \xdef\dx{13.2}\fi
  \pgfmathsetmacro{\ww}{\w*0.0365}
  \node[draw=busd, fill=busf, minimum width=\ww cm, minimum height=0.8cm, anchor=north west, inner sep=0pt, font=\tiny, align=center]
    at (\dx,10.6) {\nm\\\w};
  \pgfmathsetmacro{\nx}{\dx+\ww}\xdef\dx{\nx}}
\foreach \k in {0,1,2}{
  \node[draw=busd, fill=busf, rounded corners=3pt, minimum width=3.3cm, minimum height=1.6cm, text width=3.1cm, align=center,
        font=\tiny, anchor=north west] (dsc\k) at ({13.2+\k*4.0},9.2)
    {\textbf{descriptor \k}\\copy tile \k\\next = \ifnum\k=2 0 (end)\else descriptor \the\numexpr\k+1\relax\fi};}
\draw[dat] (dsc0) -- (dsc1); \draw[dat] (dsc1) -- (dsc2);
\node[note, font=\scriptsize, anchor=north west, text width=11.9cm] at (13.2,7.3) {%
  \textbf{The DMA engine} loads a descriptor into its registers (flip-flops), runs its address generator:\\[2pt]
  \hspace*{1em}\rn{for r in 0..rows-1: copy length bytes from src + r$\cdot$src\_stride to dst + r$\cdot$dst\_stride}\\[2pt]
  then follows \rn{next}. The CPU writes the list once and starts the DMA with one register write.\\[4pt]
  \textbf{Strides} let one descriptor copy a 2-D tile out of a big matrix, e.g. a 32$\times$32 int8 tile of a 1024-wide
  matrix: length = 32, rows = 32, src stride = 1024.\\[4pt]
  In your project \rn{host.py} plays the DMA: it sends every A and B byte itself, one bus write per byte.};

% =====================================================================
% (D) coherence problem
% =====================================================================
\panel{25.7}{0.7}{38.3}{11.9}{The cache trap: the NPU cannot see the CPU's cache}
\node[cpu, minimum width=3.4cm, minimum height=1.5cm, text width=3.2cm, anchor=north west, font=\tiny] at (26.0,11.1)
  {\textbf{CPU D-cache}\\line X = \textbf{new input}\\D = 1 (dirty)};
\node[mem, minimum width=3.4cm, minimum height=1.5cm, text width=3.2cm, anchor=north west, font=\tiny] at (30.2,11.1)
  {\textbf{DRAM}\\line X = \textbf{old data}};
\node[npu, minimum width=3.2cm, minimum height=1.5cm, text width=3.0cm, anchor=north west, font=\tiny] at (34.4,11.1)
  {\textbf{NPU DMA}\\reads X from DRAM\\$\Rightarrow$ gets old data!};
\draw[dat] (34.4,10.35) -- (33.6,10.35);
\node[font=\Large, text=bad] at (34.0,10.8) {$\times$};
\node[note, anchor=north west, text width=12.0cm] at (26.0,9.1) {%
  The CPU's write is still in its cache. DRAM has old data, and the NPU reads DRAM.\\[4pt]
  \textbf{Fix 1, software}: clean (write back) the input buffer before START (step 3); invalidate the output buffer after the
  interrupt (step 15), or the CPU may read stale cached copies.\\[4pt]
  \textbf{Fix 2, hardware}: connect the NPU through an \textbf{IO-coherent} port (ACE-Lite): its reads are checked against
  the CPU caches automatically. Costs snoop traffic.\\[4pt]
  \textbf{Fix 3}: map the buffers as uncached. Simple, but CPU access to them becomes slow.\\[4pt]
  Your project has no cache in the path, so this cannot happen today. It will matter once a CPU with a data cache
  writes the A and B buffers (Figure 12).};

\figfoot{Register names, addresses and the step order are one typical driver flow. Real drivers batch many jobs and may use several interrupts.}
\end{tikzpicture}
\end{document}
